\documentclass[11pt, a4paper]{article}

\usepackage[a4paper, margin=2.5cm, top=2.5cm, bottom=2.5cm]{geometry}
\usepackage[UTF8,fontset=mac]{ctex}
\setCJKmainfont[
    BoldFont = {Heiti SC},
    BoldFeatures = {FakeBold = 3}
]{Songti SC}

\usepackage{fontspec}
\usepackage{amsmath, amssymb, amsthm}
\usepackage{booktabs}
\usepackage{tabularx}
\usepackage{fancyhdr}
\usepackage{enumitem}
\usepackage{xcolor}
\usepackage{hyperref}
\usepackage{CJKfntef}
\usepackage{printlen} % 引入打印长度的宏包
\usepackage{array}
\usepackage{listings}    % 代码块支持
\usepackage{titlesec}
\usepackage{graphicx}     % 图片导入
\usepackage{float}        % 动态格式
\usepackage{multirow}
\usepackage{fontspec}

\setmonofont{Menlo}

% 一级标题：居中、大号、加粗
\titleformat{\section}
    {\centering\Large\bfseries}
    {\thesection}
    {1em}
    {}
% 二级标题：左对齐、较大、加粗
\titleformat{\subsection}
    {\large\bfseries}
    {}
    {0pt}
    {【}
    [】]

% 页眉页脚设置
\pagestyle{fancy}
\fancyhf{}
\newcommand{\currentproblem}{}% 当前题目的中文名称
\fancyhead[L]{\small 2026 CCF 非专业级软件能力认证 CSP-J/S 2026 第二轮} % 左上角：考试名称
\fancyhead[R]{\small 提高组（非官方试卷）\ \currentproblem} % 右上角：组别 + 当前题目名称
\fancyfoot[C]{%
    \small 第 \thepage\ 页 \hspace{1em} 共 \textcolor{blue}{\pdfpagerefbox} 页
} % 页脚：总页数用蓝色显示
\renewcommand{\headrulewidth}{0.4pt}
\renewcommand{\footrulewidth}{0pt}
\newcommand{\pdfpagerefbox}{13} % 手动设置总页数
% 首页专用页眉页脚样式
\fancypagestyle{firstpage}{
    \fancyhf{}
    % \fancyfoot[C]{%
    %     \small 第 \thepage\ 页 / 共 \textcolor{blue}{\pdfpagerefbox} 页%
    % }
    \renewcommand{\headrulewidth}{0pt}
}

% 自定义环境与命令
% 每道题作为一级 section
\newcommand{\problemheading}[2]{%
    \clearpage
    \renewcommand{\currentproblem}{#1（#2）}%
    \section*{#1 (#2)}%
}
% 题目内部标题作为二级 subsection
\newcommand{\sectionbox}[1]{
    \vspace{1.3em}%
    \noindent\hspace*{1em}%
    {\large\bfseries【#1】}\par
    \vspace{0.4em}%
}
% 题目内部标题作为第一个二级 subsection
\newcommand{\sectionboxfirst}[1]{%
    \noindent\hspace*{1em}%
    {\large\bfseries【#1】}\par
    \vspace{0.4em}%
}

% 配置代码块样式 (Consolas字体, 行号, 语法高亮)
\lstset{
    basicstyle=\ttfamily\small,
    frame=single,
    rulecolor=\color{blue},
    framerule=0.4pt,
    breaklines=true,
    keywordstyle=\color{blue},
    commentstyle=\color{green!60!black},
    stringstyle=\color{red},
    numbers=left,
    numberstyle=\footnotesize\color{gray},
    numbersep=8pt,
    xleftmargin=1.5em,
    framexleftmargin=0pt,
    tabsize=4,
    columns=fixed,
    keepspaces=true
}

% 缩短枚举元素间距离
\setlist[itemize]{
    itemsep=0pt,
    parsep=0pt,
    topsep=2pt,
    partopsep=0pt
}
\setlist[enumerate]{
    itemsep=0pt,
    parsep=0pt,
    topsep=2pt,
    partopsep=0pt
}

\setlength{\intextsep}{0pt}

\begin{document}

% ------------------ Page 1: 首页须知 ------------------
\thispagestyle{firstpage}

\begin{center}
    {\huge \textbf{2026 CCF 非专业级软件能力认证}}\\[0.5em]
    {\huge CSP-J/S 2026 \textbf{第二轮}}\\[0.5em]
    {\Large 提高组（非官方试卷）}\\[0.5em]
    {\large 时间: 2026年10月31日 14:30 $\sim$ 18:30}
\end{center}

\vspace{1em}

\noindent
\begin{tabular}{|p{3cm}|p{2.7cm}|p{2.7cm}|p{2.7cm}|p{2.7cm}|}
    \hline
    题目名称 & 能量水晶 & 地铁 & 迁徙 & 研发基地 \\ \hline
    题目类型 & 传统型 & 传统型 & 传统型 & 传统型 \\ \hline
    目录 & \texttt{crystal} & \texttt{subway} & \texttt{migrate} & \texttt{hub} \\ \hline
    可执行文件名 & \texttt{crystal} & \texttt{subway} & \texttt{migrate} & \texttt{hub} \\ \hline
    输入文件名 & \texttt{crystal.in} & \texttt{subway.in} & \texttt{migrate.in} & \texttt{hub.in} \\ \hline
    输出文件名 & \texttt{crystal.out} & \texttt{subway.out} & \texttt{migrate.out} & \texttt{hub.out} \\ \hline
    每个测试点时限 & 1.0秒 & 1.0秒 & 1.0秒 & 3.0秒 \\ \hline
    内存限制 & 1024 MiB & 1024 MiB & 1024 MiB & 1024 MiB \\ \hline
    测试点数目 & 20 & 25 & 25 & 20 \\ \hline
    测试点是否等分 & 是 & 是 & 是 & 是 \\ \hline
\end{tabular}

\vspace{1em}

提交源程序文件名

\vspace{0.5em}

\noindent
\begin{tabular}{|p{3cm}|p{2.7cm}|p{2.7cm}|p{2.7cm}|p{2.7cm}|}
    \hline
    对于 C++ 语言 & \texttt{crystal.cpp} & \texttt{subway.cpp} & \texttt{migrate.cpp} & \texttt{hub.cpp} \\ \hline
\end{tabular}

\vspace{0.5em}

编译选项

\vspace{0.5em}

\noindent
\begin{tabular}{|p{3.0cm}|>{\centering\arraybackslash}p{12.1cm}|}
    \hline
    对于 C++ 语言 & \texttt{-O2 -std=c++14 -static} \\ \hline
\end{tabular}

\vspace{1em}

\CJKunderdot{\textbf{注意事项（请仔细阅读）}}
\begin{enumerate}[label=\arabic*., itemsep=2pt, topsep=4pt]
    \item 文件名（程序名和输入输出文件名）必须使用英文小写。
    \item \texttt{main()} 函数的返回值类型必须是 \texttt{int}，程序正常结束时的返回值必须是 0。
    \item 若无特殊说明，结果的比较方式为全文比较（过滤行末空格及文末换行）。
    \item 选手提交的程序源文件大小不得超过 100KiB。
    \item 程序可使用的栈空间内存限制与题目的内存限制一致。
    \item 禁止在提交的程序源文件中改变编译器参数（如使用 \texttt{\#pragma} 命令），禁止使用系统结构相关指令（如内联汇编）或其他可能造成不公平的方法。
    \item 提交的程序源文件的放置位置请参考各省的具体要求。
    \item 因违反上述规定而出现的问题，申诉时一律不予受理。
    \item 只提供 Linux 格式附加样例文件。
    \item 全国统一评测时采用的机器配置为：Intel Core Ultra 9 285K CPU @ 3.70 GHz（关闭睿频与能效核），内存 96 GB。上述时限以此配置为准。
    \item 评测在当前最新公布的 NOI Linux 下进行，各语言的编译器版本以此为准。
\end{enumerate}

% ------------------ Page 2: 能量水晶（crystal） ------------------

\problemheading{能量水晶}{crystal}

\sectionboxfirst{题目描述}

小 R 有 $n$ 颗能量水晶，其中第 $i$ $(1 \le i \le n)$ 颗的能量值为整数 $a_i$。

小 R 想用其中若干颗能量水晶合成一颗大水晶，用于合成的水晶必须具有相同的能量值。小 R 希望参与合成的水晶尽可能多。他可以进行\CJKunderdot{\textbf{至多}}\,$m$ 次操作。每次操作可以任选一颗水晶，将其能量值增加 $1$；同一颗水晶可以被重复选择。

你需要帮助小 R 求出，进行至多 $m$ 次操作后，他最多能用多少颗能量水晶合成一颗大水晶。

\sectionbox{输入格式}

从文件 $\boldsymbol{crystal.in}$ 中读入数据。

输入的第一行包含一个正整数 $n$ 和一个非负整数 $m$，分别表示能量水晶的数量和最多可以进行的操作次数。

输入的第二行包含 $n$ 个整数 $a_1,a_2,\ldots,a_n$，依次表示每颗能量水晶的能量值。

\sectionbox{输出格式}

输出到文件 $\boldsymbol{crystal.out}$ 中。

输出一行一个正整数，表示进行至多 $m$ 次操作后，小 R 最多能用多少颗能量水晶合成一颗大水晶。

\sectionbox{样例1输入}
\begin{lstlisting}
5 5
1 2 4 8 9
\end{lstlisting}

\sectionbox{样例1输出}
\begin{lstlisting}
3
\end{lstlisting}

\sectionbox{样例1解释}
小 R 有 $5$ 颗能量水晶，能量值分别为 $1,2,4,8,9$。他可以对能量值为 $1$ 和 $2$ 的水晶分别操作 $3$ 次和 $2$ 次，使原来能量值为 $1,2,4$ 的三颗水晶的能量值都变成 $4$，共进行 $5$ 次操作。因此小 R 可以用 $3$ 颗能量水晶合成一颗大水晶。

\sectionbox{样例2输入}
\begin{lstlisting}
5 8
-3 -1 2 2 7
\end{lstlisting}

\sectionbox{样例2输出}
\begin{lstlisting}
4
\end{lstlisting}

\sectionbox{样例3}
见题目目录下的 $\boldsymbol{crystal/3.in}$ 与 $\boldsymbol{crystal/3.ans}$。

该样例满足测试点 $8,9$ 的约束条件。

\sectionbox{样例4}
见题目目录下的 $\boldsymbol{crystal/4.in}$ 与 $\boldsymbol{crystal/4.ans}$。

该样例满足测试点 $10\sim 12$ 的约束条件。

\sectionbox{数据范围}

对于所有测试数据，均有：
\begin{itemize}
    \item $1 \le n \le 5 \times 10^5$，$0 \le m \le 10^{18}$；
    \item 对于所有 $i$ $(1 \le i \le n)$，均有 $|a_i| \le 10^{12}$。
\end{itemize}

\vspace{0.5em}
\noindent
\begin{center}
\renewcommand{\arraystretch}{1.15}
\setlength{\arrayrulewidth}{0.5pt}
\begin{tabular}{c!{\vrule width 0.6pt}c!{\vrule width 0.6pt}c!{\vrule width 0.6pt}c!{\vrule width 0.6pt}c}
    \noalign{\hrule height 1.5pt}
    测试点编号 & $n\le$ & $m\le$ & $\max_{i=1}^{n} |a_i|\le$ & 特殊性质 \\
    \noalign{\hrule height 1.5pt}
    $1,2$ & $5$ & $10$ & \multirow{2}{*}{20} & \multirow{5}{*}{无} \\
    \cline{1-3}
    $3\sim5$ & $16$ & $1000$ & & \\
    \cline{1-4}
    $6,7$ & \multirow{2}{*}{5000} & $10^6$ & 100 & \\
    \cline{1-1}\cline{3-4}
    $8,9$ & & $10^9$ & $10^5$ & \\
    \cline{1-4}
    $10\sim12$ & \multirow{4}{*}{$5\times10^5$} & $10^8$ & $100$ & \\
    \cline{1-1}\cline{3-5}
    $13\sim15$ & & \multirow{3}{*}{$10^{18}$} & \multirow{3}{*}{$10^{12}$} & A \\
    \cline{1-1}\cline{5-5}
    $16\sim18$ & & & & B \\
    \cline{1-1}\cline{5-5}
    $19,20$ & & & & 无 \\
    \cline{1-1}\cline{5-5}
    \noalign{\hrule height 1.5pt}
\end{tabular}
\end{center}

特殊性质 A：对于所有 $i$ $(1 \le i \le n)$，均有 $a_i=i$。

特殊性质 B：$a_1,a_2,\ldots,a_n$ 构成等差数列。

% ------------------ Page 4: 地铁（subway） ------------------

\problemheading{地铁}{subway}

\sectionboxfirst{题目描述}

B 市共有 $m$ 个地铁站，编号为 $1$ 到 $m$。每个地铁站都有一个站厅，地铁口与站厅之间有一条通道，往返耗时相同。从第 $v$（$1 \le v \le m$）个地铁站的地铁口走到站厅需要 $w_v$ 的时间。

B 市共有 $n$ 条地铁线路，编号为 $1$ 到 $n$，其中第 $i$（$1 \le i \le n$）条线路经过 $k_i$ 个地铁站，编号依次为 $a_{i,1},a_{i,2},\ldots,a_{i,k_i}$。在第 $a_{i,j}$（$1 \le j \le k_i$）个地铁站处，从站厅走到第 $i$ 条线路的站台需要 $t_{i,j}$ 的时间，往返耗时相同。同一条线路不会重复经过同一个地铁站，即 $a_{i,1},a_{i,2},\ldots,a_{i,k_i}$ 两两不同。

由于空间有限，每个地铁站最多有 $d$ 条线路经过。此外，B 市的地铁网络是连通的，即任意两个地铁站都可以通过若干条线路互相到达。

小 R 有 $q$ 个出行计划。对于第 $r$（$1 \le r \le q$）个出行计划，小 R 将从第 $x_r$ 个地铁站的地铁口出发，乘坐地铁到达第 $y_r$ 个地铁站的地铁口，途中可能进行若干次站内换乘。每次站内换乘都需要从一条线路的站台先走到站厅，再从站厅走到另一条线路的站台。每个出行计划的总步行时间为进出站和换乘的步行时间之和。小 R 希望每个出行计划的总步行时间都尽可能少。

你需要帮助小 R 求出，每个出行计划的最小总步行时间。

\sectionbox{输入格式}

从文件 $\boldsymbol{subway.in}$ 中读入数据。

输入的第一行包含五个非负整数 $c,n,m,d,q$，分别表示测试点编号、地铁线路数、地铁站数、每个地铁站最多被几条线路经过，以及出行计划数。特别地，若 $c=0$，则该测试点为样例。

输入的第二行包含 $m$ 个正整数 $w_1,w_2,\ldots,w_m$，依次表示从每个地铁站的地铁口走到站厅需要的时间。

输入的第 $i+2$（$1 \le i \le n$）行包含 $1+2k_i$ 个正整数，其中第一个正整数为 $k_i$，表示第 $i$ 条线路经过的地铁站数；接下来 $2k_i$ 个正整数为 $a_{i,1},t_{i,1},a_{i,2},t_{i,2},\ldots,a_{i,k_i},t_{i,k_i}$，其中 $a_{i,j}$ 表示第 $i$ 条线路经过的第 $j$ 个地铁站，$t_{i,j}$ 表示从该地铁站的站厅走到第 $i$ 条线路的站台需要的时间。

输入的第 $r+n+2$（$1 \le r \le q$）行包含两个正整数 $x_r,y_r$，分别表示第 $r$ 个出行计划的起点和终点所在的地铁站编号。

\sectionbox{输出格式}

输出到文件 $\boldsymbol{subway.out}$ 中。

输出 $q$ 行，其中第 $r$（$1 \le r \le q$）行包含一个非负整数，表示第 $r$ 个出行计划的最小总步行时间。

\sectionbox{样例1输入}
\begin{lstlisting}
0 2 5 20 2
4 1 9 2 5
3 1 5 2 7 4 6
3 2 3 3 4 5 7
1 3
1 2
\end{lstlisting}

\sectionbox{样例1输出}
\begin{lstlisting}
32
17
\end{lstlisting}

\sectionbox{样例1解释}

B 市有 $5$ 个地铁站，从各地铁口走到站厅需要的时间分别为 $4,1,9,2,5$。

B 市有 $2$ 条地铁线路，其中第 $1$ 条线路依次经过第 $1,2,4$ 个地铁站，第 $2$ 条线路依次经过第 $2,3,5$ 个地铁站，两条线路在第 $2$ 个地铁站相交。

对于第 $1$ 个出行计划，小 R 将从第 $1$ 个地铁站前往第 $3$ 个地铁站：他可以先从第 $1$ 个地铁站的地铁口走到站厅，再走到第 $1$ 条线路的站台，乘坐第 $1$ 条线路到达第 $2$ 个地铁站后走回站厅，再走到第 $2$ 条线路的站台，乘坐第 $2$ 条线路到达第 $3$ 个地铁站后走回站厅，最后走到第 $3$ 个地铁站的地铁口，总步行时间为
$$
w_1+t_{1,1}+(t_{1,2}+t_{2,1})+t_{2,2}+w_3=4+5+(7+3)+4+9=32.
$$

对于第 $2$ 个出行计划，小 R 将从第 $1$ 个地铁站前往第 $2$ 个地铁站：他可以先从第 $1$ 个地铁站的地铁口走到站厅，再走到第 $1$ 条线路的站台，乘坐第 $1$ 条线路到达第 $2$ 个地铁站后走回站厅，最后走到第 $2$ 个地铁站的地铁口，总步行时间为
$$
w_1+t_{1,1}+t_{1,2}+w_2=4+5+7+1=17.
$$

\sectionbox{样例2输入}
\begin{lstlisting}
0 3 7 20 2
8 3 6 1 7 2 9
3 1 4 2 5 5 100
3 2 6 3 1 7 2
4 3 2 5 100 4 6 6 7
1 6
2 3
\end{lstlisting}

\sectionbox{样例2输出}
\begin{lstlisting}
35
16
\end{lstlisting}

\sectionbox{样例2解释}

B 市有 $7$ 个地铁站与 $3$ 条地铁线路，其中第 $1$ 条线路依次经过第 $1,2,5$ 个地铁站，第 $2$ 条线路依次经过第 $2,3,7$ 个地铁站，第 $3$ 条线路依次经过第 $3,5,4,6$ 个地铁站。

对于第 $1$ 个出行计划，小 R 将从第 $1$ 个地铁站前往第 $6$ 个地铁站：他可以先从第 $1$ 个地铁站的地铁口走到站厅，再走到第 $1$ 条线路的站台，乘坐第 $1$ 条线路到达第 $2$ 个地铁站后走回站厅，再走到第 $2$ 条线路的站台，乘坐第 $2$ 条线路到达第 $3$ 个地铁站后走回站厅，再走到第 $3$ 条线路的站台，乘坐第 $3$ 条线路到达第 $6$ 个地铁站后走回站厅，最后走到第 $6$ 个地铁站的地铁口，总步行时间为
$$
w_1+t_{1,1}+(t_{1,2}+t_{2,1})+(t_{2,2}+t_{3,1})+t_{3,4}+w_6=35.
$$

对于第 $2$ 个出行计划，小 R 将从第 $2$ 个地铁站前往第 $3$ 个地铁站：他可以先从第 $2$ 个地铁站的地铁口走到站厅，再走到第 $2$ 条线路的站台，乘坐第 $2$ 条线路到达第 $3$ 个地铁站后走回站厅，最后走到第 $3$ 个地铁站的地铁口，总步行时间为
$$
w_2+t_{2,1}+t_{2,2}+w_3=16.
$$

\sectionbox{样例3}
见题目目录下的 $\boldsymbol{subway/3.in}$ 与 $\boldsymbol{subway/3.ans}$。

该样例满足测试点 $4\sim 6$ 的约束条件。

\sectionbox{样例4}
见题目目录下的 $\boldsymbol{subway/4.in}$ 与 $\boldsymbol{subway/4.ans}$。

该样例满足测试点 $9\sim 11$ 的约束条件。

\sectionbox{样例5}
见题目目录下的 $\boldsymbol{subway/5.in}$ 与 $\boldsymbol{subway/5.ans}$。

该样例满足测试点 $12\sim 15$ 的约束条件。

\sectionbox{样例6}
见题目目录下的 $\boldsymbol{subway/6.in}$ 与 $\boldsymbol{subway/6.ans}$。

该样例满足测试点 $16\sim 19$ 的约束条件。

\sectionbox{样例7}
见题目目录下的 $\boldsymbol{subway/7.in}$ 与 $\boldsymbol{subway/7.ans}$。

该样例满足测试点 $22,23$ 的约束条件。

\clearpage % ！！！找到合适尺寸后移除

\sectionbox{数据范围}

对于所有测试数据，均有：

\begin{itemize}
    \item $0 \le c \le 25$，$1 \le n \le 500$，$1 \le m \le \sum_{i=1}^{n} k_i \le 5 \times 10^5$，$1 \le d \le 20$，$1 \le q \le 10^5$；
    \item 对于所有 $v$（$1 \le v \le m$），均有 $1 \le w_v \le 10^9$；
    \item 对于所有 $i$（$1 \le i \le n$），均有 $1 \le k_i \le m$；
    \item 对于所有 $i,j$（$1 \le i \le n$，$1 \le j \le k_i$），均有 $1 \le a_{i,j} \le m$，$1 \le t_{i,j} \le 10^9$；
    \item 对于所有 $i$（$1 \le i \le n$），均有 $a_{i,1},a_{i,2},\ldots,a_{i,k_i}$ 两两不同；
    \item 对于所有 $v$（$1 \le v \le m$），经过第 $v$ 个地铁站的线路数均不超过 $d$；
    \item 任意两个地铁站都可以通过若干条线路互相到达；
    \item 对于所有 $r$（$1 \le r \le q$），均有 $1 \le x_r,y_r \le m$，且 $x_r \ne y_r$。
\end{itemize}

\vspace{0.5em}
\noindent
\begin{center}
\renewcommand{\arraystretch}{1.15}
\setlength{\arrayrulewidth}{0.5pt}
\begin{tabular}{c!{\vrule width 0.6pt}c!{\vrule width 0.6pt}c!{\vrule width 0.6pt}c!{\vrule width 0.6pt}c!{\vrule width 0.6pt}c}
    \noalign{\hrule height 1.5pt}
    测试点编号 & $n\le$ & $\sum_{i=1}^{n} k_i\le$ & $d\le$ & $q\le$ & 特殊性质 \\
    \noalign{\hrule height 1.5pt}
    $1\sim3$ & 3 & 10 & \multirow{7}{*}{20} & 10 & \multirow{2}{*}{无} \\
    \cline{1-3}\cline{5-5}
    $4\sim6$ & 50 & 300 & & \multirow{2}{*}{200} & \\
    \cline{1-3}\cline{6-6}
    $7,8$ & 200 & 400 & & & \multirow{2}{*}{A} \\
    \cline{1-3}\cline{5-5}
    $9\sim11$ & \multirow{3}{*}{500} & 1000 & & $10^5$ & \\
    \cline{1-1}\cline{3-3}\cline{5-6}
    $12\sim15$ & & $2\times10^4$ & & 2000 & \multirow{2}{*}{B} \\
    \cline{1-1}\cline{3-3}\cline{5-5}
    $16\sim19$ & & \multirow{4}{*}{$5\times10^5$} & & \multirow{4}{*}{$10^5$} & \\
    \cline{1-2}\cline{6-6}
    $20,21$ & 2 & & & & \multirow{3}{*}{无} \\
    \cline{1-2}\cline{4-4}
    $22,23$ & \multirow{2}{*}{500} & & 2 & & \\
    \cline{1-1}\cline{4-4}
    $24,25$ & & & 20 & & \\
    \noalign{\hrule height 1.5pt}
\end{tabular}
\end{center}

特殊性质 A：每条线路至多经过 $2$ 个地铁站，即对于所有 $i$（$1 \le i \le n$），均有 $k_i \le 2$。

特殊性质 B：对于所有线路经过的所有地铁站，从站厅走到对应线路站台的步行时间都相同，即存在正整数 $T$，使得对于所有 $i,j$（$1 \le i \le n$，$1 \le j \le k_i$），均有 $t_{i,j}=T$。

% ------------------ Page 8: 迁徙（migrate） ------------------

\problemheading{迁徙}{migrate}

\sectionboxfirst{题目描述}

海上有 $n$ 座岛屿，自北向南排成一条直线，其中第 $i$（$1 \le i \le n$）座岛屿在这条直线上的坐标为整数 $x_i$，且 $x_1 < x_2 < \cdots < x_n$。每年秋季，一群海鸟会沿着这条直线在南北方向上迁徙，小 P 正在观察它们。

海鸟的体力有限，每次飞行的距离都不超过 $d$。若两座岛屿 $i,j$（$1 \le i,j \le n$）满足 $|x_i-x_j| \le d$，海鸟就可以从第 $i$ 座岛屿直接飞到第 $j$ 座岛屿。

对于两座不同的岛屿 $u,v$（$1 \le u,v \le n$），一条从第 $u$ 座岛屿到第 $v$ 座岛屿的路径形如 $y_1 \to y_2 \to \cdots \to y_m$，其中 $y_1=u$，$y_m=v$，且对于所有 $l$（$1 \le l \le m-1$），海鸟都能够从第 $y_l$ 座岛屿直接飞到第 $y_{l+1}$ 座岛屿；其长度为 $m-1$，即海鸟在路径中飞行的次数。

小 P 发现，海鸟在迁徙时总是会选择长度最短的路径。对于两座岛屿 $u,v$（$1 \le u < v \le n$），若从第 $u$ 座岛屿到第 $v$ 座岛屿的长度最短的路径\CJKunderdot{\textbf{恰好}}有两条，海鸟可以分成两群，分别沿其中一条飞往第 $v$ 座岛屿；此时称 $(u,v)$ 为\CJKunderdot{\textbf{双径岛屿对}}。两条路径不同，当且仅当它们经过的岛屿序列不完全相同。

你需要帮助小 P 求出双径岛屿对的数量。

\sectionbox{输入格式}

从文件 $\boldsymbol{migrate.in}$ 中读入数据。

\CJKunderdot{\textbf{本题包含多组测试数据。}}

输入的第一行包含两个非负整数 $c,t$，分别表示测试点编号与测试数据组数。特别地，若 $c=0$，则该测试点为样例。

接下来依次输入 $t$ 组测试数据。对于每组测试数据：

\begin{itemize}
    \item 第一行包含两个正整数 $n,d$，分别表示岛屿数量和海鸟一次最多能飞行的距离。
    \item 第二行包含 $n$ 个严格递增的整数 $x_1,x_2,\ldots,x_n$，依次表示每座岛屿的坐标。
\end{itemize}

\sectionbox{输出格式}

输出到文件 $\boldsymbol{migrate.out}$ 中。

对于每组测试数据，输出一行一个非负整数，表示双径岛屿对的数量。

\sectionbox{样例1输入}
\begin{lstlisting}
0 2
4 2
1 2 3 4
5 2
1 2 3 4 5
\end{lstlisting}

\sectionbox{样例1输出}
\begin{lstlisting}
1
2
\end{lstlisting}

\sectionbox{样例1解释}

该样例共包含两组测试数据。

对于第一组测试数据，仅有 $1$ 个双径岛屿对 $(1,4)$。从第 $1$ 座岛屿到第 $4$ 座岛屿有两条长度最短的路径，分别为 $1 \to 2 \to 4$ 和 $1 \to 3 \to 4$。

对于第二组测试数据，共有 $2$ 个双径岛屿对，分别为 $(1,4)$ 和 $(2,5)$。从第 $1$ 座岛屿到第 $5$ 座岛屿只有一条长度最短的路径 $1 \to 3 \to 5$，因此 $(1,5)$ 不是双径岛屿对。

\sectionbox{样例2输入}
\begin{lstlisting}
0 2
5 3
-2 -1 0 1 2
6 2
-4 -3 -2 -1 3 6
\end{lstlisting}

\sectionbox{样例2输出}
\begin{lstlisting}
0
1
\end{lstlisting}

\sectionbox{样例2解释}

该样例共包含两组测试数据。

对于第一组测试数据，共有 $0$ 个双径岛屿对。从第 $1$ 座岛屿到第 $5$ 座岛屿有三条长度最短的路径，分别为 $1 \to 2 \to 5$，$1 \to 3 \to 5$ 和 $1 \to 4 \to 5$，因此 $(1,5)$ 不是双径岛屿对。

对于第二组测试数据，仅有 $1$ 个双径岛屿对 $(1,4)$。从第 $1$ 座岛屿到第 $4$ 座岛屿有两条长度最短的路径，分别为 $1 \to 2 \to 4$ 和 $1 \to 3 \to 4$。第 $5$ 座和第 $6$ 座岛屿均无法到达任何其他岛屿。

\sectionbox{样例3}
见题目目录下的 $\boldsymbol{migrate/3.in}$ 与 $\boldsymbol{migrate/3.ans}$。

该样例满足测试点 $3\sim 6$ 的约束条件。

\sectionbox{样例4}
见题目目录下的 $\boldsymbol{migrate/4.in}$ 与 $\boldsymbol{migrate/4.ans}$。

该样例满足测试点 $11,12$ 的约束条件。

\sectionbox{样例5}
见题目目录下的 $\boldsymbol{migrate/5.in}$ 与 $\boldsymbol{migrate/5.ans}$。

该样例满足测试点 $15\sim18$ 的约束条件。

\sectionbox{样例6}
见题目目录下的 $\boldsymbol{migrate/6.in}$ 与 $\boldsymbol{migrate/6.ans}$。

该样例满足测试点 $19\sim21$ 的约束条件。

\sectionbox{数据范围}

设 $N$ 为单个测试点内所有测试数据的 $n$ 的和。对于所有测试数据，均有：

\begin{itemize}
    \item $0 \le c \le 25$，$2 \le t \le 10^5$；
    \item $1 \le n \le 2 \times 10^6$，$N \le 2 \times 10^6$，$1 \le d \le 2 \times 10^{18}$；
    \item $-10^{18} \le x_1 < x_2 < \cdots < x_n \le 10^{18}$。
\end{itemize}

\vspace{0.5em}
\noindent
\begin{center}
\renewcommand{\arraystretch}{1.15}
\setlength{\arrayrulewidth}{0.5pt}
\begin{tabular}{c!{\vrule width 0.6pt}c!{\vrule width 0.6pt}c}
    \noalign{\hrule height 1.5pt}
    测试点编号 & $N\le$ & 特殊性质\\
    \noalign{\hrule height 1.5pt}
    $1,2$ & 9 & \multirow{3}{*}{无} \\
    \cline{1-2}
    $3\sim6$ & 300 & \\
    \cline{1-2}
    $7\sim10$ & 5000 & \\
    \cline{1-3}
    $11,12$ & \multirow{2}{*}{$2\times10^6$} & A \\
    \cline{1-1}\cline{3-3}
    $13,14$ & & B \\
    \cline{1-3}
    $15\sim18$ & $10^6$ & C \\
    \cline{1-3}
    $19\sim21$ & $2\times10^5$ & \multirow{2}{*}{无} \\
    \cline{1-2}
    $12\sim25$ & $2\times10^6$ & \\
    \noalign{\hrule height 1.5pt}
\end{tabular}
\end{center}

特殊性质 A：$x_1,x_2,\ldots,x_n$ 构成等差数列。

特殊性质 B：存在一座岛屿，使得海鸟可以从该岛屿直接飞到任意其他岛屿。

特殊性质 C：$x_n-x_1 \le 10d$。

% ------------------ Page 11: 研发基地（hub） ------------------

\problemheading{研发基地}{hub}

\sectionboxfirst{题目描述}

小 W 曾是一名科研爱好者。在 C 国工作时，他担任测试员，日复一日地测试各种新型传送门，在城市之间来回穿梭。久而久之，他厌倦了这样的生活，于是辞去工作，来到没有任何传送门的 D 国，希望在这里重新找回对研发的热爱。

在开始研发之前，小 W 需要建造研发基地。D 国共有 $n$ 座城市，编号为 $1$ 到 $n$。这 $n$ 座城市由 $n-1$ 条道路连接，其中第 $i\ (1\le i\le n-1)$ 条道路连接第 $u_i$ 座和第 $v_i$ 座城市。若将每座城市视作一个点，将每条道路视作连接两个点的边，则这 $n$ 座城市和 $n-1$ 条道路构成一棵树 $T$。

小 W 的经费有限，他计划选择\CJKunderdot{\textbf{恰好}}\,$k$ 座城市，在每座选中的城市里各建造一个研发基地。为了便于这些基地之间的协作，小 W 希望选中的城市在 $T$ 中对应的点形成\CJKunderdot{\textbf{不超过两个连通块}}。形式化地，定义建造方案 $S\subseteq\{1,2,\ldots,n\}$ 是\CJKunderdot{\textbf{合理的}}，当且仅当 $|S|=k$，且 $S$ 在 $T$ 中形成的连通块数量不超过 $2$。

对于合理的建造方案 $S$，定义第 $i\ (1\le i\le n)$ 座城市的\CJKunderdot{\textbf{通勤距离}}为 $d(i)=\min_{j\in S}\operatorname{dist}(i,j)$，其中 $\operatorname{dist}(i,j)$ 表示从第 $i$ 座城市到第 $j$ 座城市所需经过的道路的最少条数。特别地，若 $i\in S$，则 $d(i)=0$。定义该建造方案的\CJKunderdot{\textbf{通勤半径}}为所有城市通勤距离的最大值，即 $\max_{i=1}^{n}d(i)$。

你需要帮助小 W 求出，所有合理的建造方案的通勤半径的最小值。

\sectionbox{输入格式}

从文件 $\boldsymbol{hub.in}$ 中读入数据。

\CJKunderdot{\textbf{本题包含多组测试数据。}}

输入的第一行包含两个非负整数 $c,t$，分别表示测试点编号与测试数据组数。特别地，若 $c=0$，则该测试点为样例。

接下来依次输入 $t$ 组测试数据。对于每组测试数据：

\begin{itemize}
    \item 第一行包含两个正整数 $n,k$，分别表示城市数量和小 W 计划建造的研发基地数量。
    \item 第 $i+1\ (1\le i\le n-1)$ 行包含两个正整数 $u_i,v_i$，分别表示第 $i$ 条道路连接的两座城市的编号。
\end{itemize}

\sectionbox{输出格式}

输出到文件 $\boldsymbol{hub.out}$ 中。

对于每组测试数据，输出一行一个非负整数，表示所有合理的建造方案的通勤半径的最小值。

\sectionbox{样例1输入}
\begin{lstlisting}
0 2
5 2
1 2
2 3
3 4
4 5
10 3
1 2
1 3
2 4
2 5
3 6
3 7
3 8
7 9
9 10
\end{lstlisting}

\sectionbox{样例1输出}
\begin{lstlisting}
1
2
\end{lstlisting}

\sectionbox{样例1解释}

该样例共包含两组测试数据。

对于第一组测试数据，共有 $10$ 种合理的建造方案。按通勤半径分类如下：

\begin{itemize}
    \item 通勤半径为 $1$ 的方案：$S=\{1,4\}$，$S=\{2,4\}$，$S=\{2,5\}$；
    \item 通勤半径为 $2$ 的方案：$S=\{1,3\}$，$S=\{1,5\}$，$S=\{2,3\}$，$S=\{3,4\}$，$S=\{3,5\}$；
    \item 通勤半径为 $3$ 的方案：$S=\{1,2\}$，$S=\{4,5\}$。
\end{itemize}

因此，所有合理的建造方案的通勤半径的最小值为 $1$。

对于第二组测试数据，一种合理的建造方案为 $S=\{2,7,9\}$，其中第 $2$ 座城市单独构成一个连通块，第 $7$ 座和第 $9$ 座城市构成一个连通块。该建造方案的通勤半径为 $2$。

\sectionbox{样例2}
见题目目录下的 $\boldsymbol{hub/2.in}$ 与 $\boldsymbol{hub/2.ans}$。

该样例满足测试点 $2,3$ 的约束条件。

\sectionbox{样例3}
见题目目录下的 $\boldsymbol{hub/3.in}$ 与 $\boldsymbol{hub/3.ans}$。

该样例满足测试点 $8$ 的约束条件。

\sectionbox{样例4}
见题目目录下的 $\boldsymbol{hub/4.in}$ 与 $\boldsymbol{hub/4.ans}$。

该样例满足测试点 $11,12$ 的约束条件。

\sectionbox{样例5}
见题目目录下的 $\boldsymbol{hub/5.in}$ 与 $\boldsymbol{hub/5.ans}$。

该样例满足测试点 $13,14$ 的约束条件。

\sectionbox{数据范围}

设 $N$ 为单个测试点内所有测试数据的 $n$ 的和。对于所有测试数据，均有：

\begin{itemize}
    \item $0\le c\le 20$，$1\le t\le 5$；
    \item $2\le n\le 10^6$，$N\le 2\times 10^6$，$2\le k\le n$；
    \item 对于所有 $i\ (1\le i\le n-1)$，均有 $1\le u_i,v_i\le n$，且所有 $(u_i,v_i)$ 构成一棵树。
\end{itemize}

\vspace{0.5em}
\noindent
\begin{center}
\renewcommand{\arraystretch}{1.15}
\setlength{\arrayrulewidth}{0.5pt}
\begin{tabular}{c!{\vrule width 0.6pt}c!{\vrule width 0.6pt}c!{\vrule width 0.6pt}c!{\vrule width 0.6pt}c}
    \noalign{\hrule height 1.5pt}
    测试点编号 & $n\le$ & $N\le$ & $k\le$ & 特殊性质 \\
    \noalign{\hrule height 1.5pt}
    1 & 5 & 15 & 5 & \multirow{7}{*}{无} \\
    \cline{1-4}
    $2,3$ & 10 & 30 & 10 & \\
    \cline{1-4}
    $4\sim6$ & 100 & 300 & 100 & \\
    \cline{1-4}
    7 & \multirow{2}{*}{1000} & \multirow{2}{*}{3000} & 2 & \\
    \cline{1-1}\cline{4-4}
    8 & & & 1000 & \\
    \cline{1-4}
    $9,10$ & 4000 & $10^4$ & 4000 & \\
    \cline{1-4}
    $11,12$ & \multirow{2}{*}{$10^5$} & \multirow{2}{*}{$3\times10^5$} & 2 & \\
    \cline{1-1}\cline{4-5}
    $13,14$ & & & $10^5$ & A \\
    \cline{1-5}
    15 & \multirow{2}{*}{$5\times10^5$} & \multirow{2}{*}{$10^6$} & 10 & \multirow{3}{*}{无} \\
    \cline{1-1}\cline{4-4}
    $16,17$ & & & $5\times10^5$ & \\
    \cline{1-4}
    $18\sim20$ & $10^6$ & $2\times10^6$ & $10^6$ & \\
    \noalign{\hrule height 1.5pt}  %11
\end{tabular}
\end{center}

特殊性质 A：存在正整数 $m$ 满足 $n=2m+1$，且对于所有 $i\ (1\le i\le m)$，均有 $u_i=i$ 且 $v_i=i+1$；对于所有 $i\ (m+1\le i\le n-1)$，均有 $u_i=i-m$ 且 $v_i=i+1$。

\end{document}